\section{Perbandingan Game Engine}

Terdapat berbagai game engine populer. Pemilihan engine bergantung pada jenis game, tim, dan target platform. Berikut perbandingan singkat \cite{unityManual}:

\begin{table}[htbp]
\centering
\small
\begin{tabular}{|l|>{\raggedright\arraybackslash}p{3.2cm}|>{\raggedright\arraybackslash}p{3cm}|>{\raggedright\arraybackslash}p{2.2cm}|}
\hline
\textbf{Engine} & \textbf{Keunggulan} & \textbf{Kekurangan} & \textbf{Platform} \\
\hline
Unity & Mudah, gratis, multiplatform & Build size besar & PC, Mobile, Web \\
\hline
Unreal & Grafis tinggi, Blueprint & Kurva belajar curam & PC, Console \\
\hline
Godot & Open source, ringan & Ekosistem lebih kecil & PC, Mobile \\
\hline
\end{tabular}
\caption{Perbandingan Game Engine Populer}
\end{table}

\subsection{Unreal Engine}

Unreal Engine menggunakan C++ dan sistem Blueprint (visual scripting). Cocok untuk game AAA dengan grafis film. Kurva belajar cukup curam untuk pemula. Gratis dengan royalti setelah pendapatan tertentu.

\subsection{Godot}

Godot adalah engine open source dengan GDScript (mirip Python) atau C\#. Ringan, cocok untuk game 2D dan 3D sederhana. Komunitas lebih kecil dibanding Unity.

\subsection{Mengapa Unity untuk Buku Ini?}

Unity dipilih karena kemudahan untuk pemula dan kesesuaian dengan kurikulum:

\begin{enumerate}
  \item \textbf{Tutorial Resmi}: Unity Learn menyediakan Roll a Ball dan tutorial lain yang terstruktur langkah demi langkah
  \item \textbf{Asset Store}: Banyak aset gratis untuk pembelajaran
  \item \textbf{C\#}: Bahasa yang lebih mudah daripada C++ (Unreal), cocok untuk mahasiswa
  \item \textbf{Komunitas}: Forum dan dokumentasi sangat lengkap
  \item \textbf{Multiplatform}: Satu project bisa di-build untuk Windows, Mac, Android, iOS, Web
\end{enumerate}
